\documentclass[11pt]{article}

\usepackage[margin=1in]{geometry}
\usepackage{amsmath, amssymb}
\usepackage{enumitem}
\usepackage{actuarialangle}
\usepackage{setspace}

\setlength{\parindent}{0pt}
\setlength{\parskip}{8pt}

\begin{document}

\begin{center}
    {\Large \textbf{Loans Practice}}\\
\end{center}

\vspace{10pt}

\hrule

\medskip

\textbf{Problem 1.} A purchase is financed by five annual payments of \$1{,}000, with the first payment due at the time of purchase. The effective rate of discount is \(d = 10\%\).

Find the total amount of interest paid.

\medskip
\hrule


\medskip



\medskip
\hrule

\hrule


\textbf{Problem 3.} A borrower planned to repay a loan of \(L\) with level payments at the end of each month for 30 years. The loan had an annual nominal interest rate of \(6\%\), convertible monthly. Starting with the 181st payment, the borrower increased the monthly payment to \(2{,}000\), which enabled the borrower to pay off the loan five years earlier than planned.

Calculate \(L\).

\medskip
\hrule




\textbf{Problem 4.} A loan is being repaid in five annual payments. The first two payments are \$200. The third and fourth payments are \$400. The final payment is \$500. The annual effective interest rate is \(6\%\). Determine the interest portion of the third payment.

\medskip
\hrule

\medskip

\textbf{Problem 5.} Angela borrowed \$30{,}000 for her college tuition at \(3\%\) per year convertible quarterly.

Her father agreed to make payments until the outstanding balance was less than \$10{,}000.

At that time Angela would make payments until the loan would be repaid.

The length of the loan is for 15 years and level payments are made on an annual basis.

How many payments will Angela need to make?

\newpage

%%% NEW 3 problems I, II, and III %%%%


\textbf{Problem I.} Seth repays a 30-year loan with a payment at the end of each year. Each of the first 20 payments is 1200, and each of the last 10 payments is 900. Interest on the loan is at an annual effective rate of \(i\), \(i>0\). The interest portion of the 11th payment is twice the interest portion of the 21st payment.

Calculate the interest portion of the 21st payment.

\begin{enumerate}[label=\Alph*.]
    \item 250
    \item 275
    \item 300
    \item 325
    \item There is not enough information to calculate the interest portion of the 21st payment.
\end{enumerate}

\medskip
\hrule

\medskip

\textbf{Problem II.} On December 31, 1984, Smith borrowed \$5{,}000 to be repaid in four years with level payments made at the end of every quarter. The first payment was made on March 31, 1985. The effective annual interest rate was \(4\%\). What was the amount of interest paid in 1986?

\begin{enumerate}[label=\Alph*., leftmargin=2.2em]
    \item Less than \$130
    \item At least \$130, but less than \$133
    \item At least \$133, but less than \$136
    \item At least \$136, but less than \$139
    \item \$139 or more
\end{enumerate}

\medskip
\hrule

\medskip

\textbf{Problem III.} A loan of \$5{,}000 is repaid with equal quarterly principal repayments and interest on the outstanding balance.

The interest rate charged on the loan is an annual effective rate of \(21.55\%\).

The debt is repaid with 20 quarterly payments. The first payment is made at the end of the first quarter.

Calculate the total payment in the first year.

\begin{enumerate}[label=\Alph*., leftmargin=2.2em]
    \item 1658
    \item 1725
    \item 1808
    \item 1858
    \item 1925
\end{enumerate}

\medskip
\hrule


\newpage

\textbf{Problem I.} A homeowner borrows \$1{,}000 to be repaid with payments at the end of each year for 20 years. There are two repayment options. The first option is equal annual payments based on an annual effective interest rate of \(3\%\). The second option is payments of \$50 each year plus interest on the unpaid balance at an annual effective interest rate of \(i\). The total payments under the two options are the same.

Calculate \(i\).

\begin{enumerate}[label=\Alph*.]
    \item 2.86\%
    \item 3.00\%
    \item 3.28\%
    \item 3.44\%
    \item 4.76\%
\end{enumerate}

\medskip
\hrule

\medskip

\textbf{Problem II.} A borrower takes out a loan of \$4{,}000 at an annual effective interest rate of \(6\%\).

Starting at the end of the fifth year, the loan is repaid by annual payments, each of which equals \$600 except for a final balloon payment that is less than \$1{,}000.

Calculate the final balloon payment.

\begin{enumerate}[label=\Alph*., leftmargin=2.2em]
    \item 616
    \item 639
    \item 642
    \item 688
    \item 696
\end{enumerate}

\medskip
\hrule

\medskip

\textbf{Problem III.} Nicole borrows \$50{,}000 at an annual effective interest rate of \(7\%\). She repays the loan with \(n\) annual payments at the end of each year. The first twenty payments are level. Starting with the 21st payment, each subsequent payment is \(2\%\) less than the previous payment for the remaining loan repayment schedule. The principal repaid in the twentieth payment, \(P_{20}\), is \$2{,}116.59. If \(I_t\) is the interest paid in the \(t\)-th period, find \(I_{21} + I_{22} + I_{23}\).

\begin{enumerate}[label=\Alph*.]
    \item 3{,}000
    \item 4{,}000
    \item 5{,}000
    \item 6{,}000
    \item 7{,}000
\end{enumerate}

\medskip
\hrule



\end{document}